% FPR — Topic B: Noise-Robust Test-Time Adaptation
% CVPR-format 5-page final project report
% =========================================================
% OVERLEAF: upload main.tex + references.bib, compiler = pdfLaTeX
% For official CVPR style: add cvpr.sty from the CVPR 2024 author
% kit and uncomment \usepackage[pagenumbers]{cvpr} below.
% =========================================================
\documentclass[10pt,twocolumn,letterpaper]{article}

% \usepackage[pagenumbers]{cvpr}   % uncomment if cvpr.sty is uploaded

\usepackage[letterpaper,
  left=0.875in,right=0.875in,top=1.0in,bottom=1.125in]{geometry}
\setlength{\columnsep}{0.375in}

\usepackage{times}
\usepackage{graphicx}
\usepackage{amsmath,amssymb}
\usepackage{booktabs}
\usepackage{multirow}
\usepackage{microtype}
\usepackage[font=small,labelfont=bf,skip=4pt]{caption}
\usepackage[hidelinks]{hyperref}

\usepackage{titlesec}
\titlespacing*{\section}   {0pt}{7pt plus 2pt minus 1pt}{3pt plus 1pt}
\titlespacing*{\subsection}{0pt}{5pt plus 1pt minus 1pt}{2pt plus 1pt}
\titlespacing*{\paragraph} {0pt}{4pt plus 1pt}{0.5em}

\setlength{\floatsep}{6pt plus 2pt minus 2pt}
\setlength{\textfloatsep}{8pt plus 2pt minus 2pt}
\setlength{\intextsep}{6pt plus 2pt minus 2pt}

\newcommand{\eg}{\textit{e.g.}}
\newcommand{\ie}{\textit{i.e.}}
\newcommand{\best}[1]{\textbf{#1}}
% =========================================================
\begin{document}

\title{%
  Noise-Robust Test-Time Adaptation of Vision-Language Models\\[-1pt]
  via Streaming Stability Controls\\[2pt]
  {\large Final Project Report FP17 group --- Topic D}%
}
\author{\textbf{Abzal Nurgazy} \\ \small April 26, 2026}
\date{}
\maketitle
\thispagestyle{empty}

% =========================================================
\begin{abstract}
Test-time adaptation (TTA) enables CLIP-style vision-language models
to improve accuracy under distribution shift without retraining.
TDA~\cite{karmanov2024tda}, the state-of-the-art training-free TTA
method, builds dynamic positive and negative key-value caches from the
test stream, but provides no quality gate on cache updates: a
confidently-wrong prediction is stored and corrupts subsequent logits.
We implement TDA from scratch, verify it against published results,
and extend it with three independently ablatable noise-robustness
controls --- margin filtering, momentum blending, and affinity decay
--- together with an uncertainty-aware adaptation strength as a
stretch goal.
All four controls are evaluated on OOD (ImageNet-A) and cross-domain
(DTD) benchmarks with two CLIP backbones (RN50 and ViT-B/16).
Systematic ablations over confidence threshold, shot capacity,
momentum, and decay factor are reported alongside per-image latency
to characterise the accuracy--compute trade-off of each design choice.
\end{abstract}

% =========================================================
\section{Introduction}
\label{sec:intro}

Vision-language models such as CLIP~\cite{radford2021clip} achieve
impressive zero-shot classification by jointly embedding images and
text in a shared feature space.
Classification is performed by matching an image feature against a
set of class-name text embeddings, requiring no task-specific
training data.
However, accuracy degrades significantly when the test distribution
differs from the pre-training distribution --- a problem known as
\emph{distribution shift}, which manifests as changes in image style,
texture, rendering, or adversarial statistics.

\emph{Test-time adaptation} (TTA) addresses this by adapting the
model at inference time using only unlabelled test images, without
accessing source training data or labels.
Prior TTA methods for vision-language models such as
TPT~\cite{shu2022tpt} and DiffTPT~\cite{feng2023difftpt} achieve
strong accuracy by learning a per-sample prompt via backpropagation,
but this process is computationally intensive: TPT requires 12 hours
and DiffTPT over 34 hours to process the 50,000-image ImageNet
validation set.

TDA~\cite{karmanov2024tda} avoids backpropagation entirely by
maintaining two lightweight key-value caches that accumulate
knowledge from the test stream.
A \emph{positive cache} stores high-confidence pseudo-labelled
image features to refine future predictions toward seen classes.
A \emph{negative cache} stores uncertain predictions to suppress
noisy pseudo-labels.
TDA matches or exceeds TPT accuracy while reducing testing time from
hours to 16 minutes, making it highly practical.

\paragraph{Problem.}
Despite its efficiency, TDA's cache update rule has a critical
limitation: it uses prediction entropy as the only quality criterion.
Any sample whose entropy falls below a fixed threshold is stored
regardless of how ambiguous the prediction is.
A \emph{confidently-wrong} prediction --- high softmax probability
on the wrong class --- passes the entropy gate and is stored,
shifting logits for all subsequent images of that class in the
wrong direction.
Furthermore, entries written early in the stream persist at full
weight throughout the entire test pass, even if later, higher-quality
examples arrive.
Neither of these failure modes is addressed in the original TDA paper.

\paragraph{Contributions.}
This report makes four contributions.
First, we faithfully re-implement TDA and verify our reproduction
against published results within $\pm0.4\%$ on both OOD and
cross-domain benchmarks with two CLIP backbones (Section~\ref{sec:exp}).
Second, we propose and ablate three streaming stability controls:
a margin filter, momentum blending, and affinity decay
(Section~\ref{sec:method}).
Third, we provide systematic ablations on all key hyperparameters
--- confidence threshold, shot capacity, momentum, and decay factor
--- with accuracy \emph{and} per-image latency reported for each
configuration (Section~\ref{subsec:ablations}).
Fourth, as a stretch goal, we introduce uncertainty-aware adaptation
strength and sweep its sensitivity parameter
(Section~\ref{subsec:unc}).

% =========================================================
\section{Related Work}
\label{sec:related}

\paragraph{Test-time adaptation.}
TENT~\cite{wang2021tent} performs fully test-time adaptation by
minimising prediction entropy through updates to batch-normalisation
parameters.
MEMO~\cite{zhang2022memo} achieves robustness by enforcing consistent
predictions across diverse augmentations of each single test image.
Both methods require gradient computation per test sample, which
limits their throughput.
Our work, following TDA, avoids any gradient computation at test time.

\paragraph{TTA for vision-language models.}
TPT~\cite{shu2022tpt} adapts CLIP to each test image by minimising
marginal entropy over multiple augmented views, optimising a
learnable prompt via backpropagation.
DiffTPT~\cite{feng2023difftpt} improves upon TPT by using a
pre-trained diffusion model to generate more diverse augmentations.
While effective, both methods are computationally expensive.
TDA~\cite{karmanov2024tda} replaces prompt optimisation with a
non-parametric dynamic cache, achieving comparable accuracy at
dramatically lower latency.

\paragraph{Cache-based adaptation.}
Tip-Adapter~\cite{zhang2022tipadapter} builds a static key-value
cache from a small set of labelled few-shot examples to augment
CLIP predictions.
TDA extends this idea to the test-time setting by constructing the
cache on-the-fly from unlabelled pseudo-labelled test samples.
Our controls directly address limitations of this dynamic cache
construction: noisy entries that corrupt the positive cache, and
stale entries that persist beyond their useful lifetime.

% =========================================================
\section{Method}
\label{sec:method}

\subsection{TDA Background}
\label{subsec:tda}

Let $E_v$ and $E_t$ denote the CLIP image and text encoders,
and let $\mathbf{W}_c \in \mathbb{R}^{d \times K}$ be the matrix
of $\ell_2$-normalised text embeddings for $K$ classes.
Given test image $x$, the zero-shot CLIP logits are
$\mathbf{l}_\text{zs} = 100\,E_v(x)\,\mathbf{W}_c^\top$,
where the factor 100 scales cosine similarities into a practical
logit range, and the prediction is
$\mathbf{p}_\text{zs} = \text{softmax}(\mathbf{l}_\text{zs})$.

\textbf{Positive cache.}
A per-class dynamic queue stores at most $k$ key-value pairs
$(\mathbf{f}, \hat{l})$, where $\mathbf{f} = E_v(x)$ is the
image feature and $\hat{l}$ is a one-hot pseudo-label.
When the queue is not full, new samples are appended.
When full, TDA replaces the highest-entropy entry if the incoming
sample has strictly lower entropy.
This priority-queue policy ensures the cache retains the most
confident samples seen so far.
The positive cache logit contribution is:
\begin{equation}
  \mathbf{p}_\text{pos} =
    \alpha \exp\!\bigl(-\beta(1 - \mathbf{f}\,Q_p^\top)\bigr)\,L_p,
\end{equation}
where $Q_p$ are the cached key features, $L_p$ are the one-hot
labels, and $\alpha, \beta$ are a residual ratio and sharpness factor.

\textbf{Negative cache.}
Samples with intermediate entropy
$\tau_l < H(\mathbf{p}_\text{zs}) < \tau_h$ are stored in a
negative cache using the same capacity management policy.
Their pseudo-labels are inverted: classes with probability above
a threshold $p_l$ are marked as absent.
The negative prediction $\mathbf{p}_\text{neg}$ is subtracted from
the final logits.
The full TDA prediction combines all three terms:
\begin{equation}
  \mathbf{p}_\text{TDA} =
    \mathbf{f}\,\mathbf{W}_c^\top +
    \mathbf{p}_\text{pos} +
    \mathbf{p}_\text{neg}.
\end{equation}

\subsection{Control 1 --- Margin Filter}
\label{subsec:margin}

The prediction \emph{margin} is the gap between the top-1 and top-2
logits: $m(x) = z_{(1)} - z_{(2)}$.
A small margin indicates that the model is nearly equally confident
between two classes, and such predictions are likely to be noisy or
wrong.
We gate the positive cache update: if the margin falls below a
threshold, the update is skipped entirely:
\begin{equation}
  \text{skip update if } z_{(1)} - z_{(2)} < \tau_m.
\end{equation}
Since the logits are already computed before the cache update, this
gate adds zero computational overhead.
The default threshold $\tau_m = 5.0$ corresponds to a cosine gap of
$0.05$ (logits are scaled by $100\times$ cosine similarity).
A higher threshold filters more aggressively, keeping only the most
unambiguous predictions in the cache.

\subsection{Control 2 --- Momentum Blending}
\label{subsec:momentum}

The standard TDA hard-replace policy can produce sharp, discontinuous
changes to a stored prototype when a single outlier sample arrives.
Momentum blending replaces the hard write with a smooth interpolation
between the incoming feature and the best currently-stored entry for
that class:
\begin{equation}
  \tilde{\mathbf{f}} =
    \frac{
      \mu\,\mathbf{f}_\text{old} + (1-\mu)\,\mathbf{f}_\text{new}
    }{
      \|\mu\,\mathbf{f}_\text{old} + (1-\mu)\,\mathbf{f}_\text{new}\|_2
    }.
\end{equation}
The re-normalisation keeps $\tilde{\mathbf{f}}$ on the unit sphere,
preserving compatibility with cosine similarity retrieval.
Setting $\mu = 0$ recovers the original TDA hard-replace;
$\mu \to 1$ freezes the prototype entirely.
Intermediate values ($\mu = 0.9$) allow slow adaptation while
dampening the influence of individual outliers.

\subsection{Control 3 --- Affinity Decay}
\label{subsec:decay}

Entries written early in the test stream may become stale as the
cache accumulates better, more representative prototypes.
We apply per-entry exponential decay at retrieval time:
\begin{equation}
  \hat{A}_{ij} =
    (\mathbf{f}_i \cdot \mathbf{q}_j) \cdot \gamma^{\,\text{age}_j},
\end{equation}
where $\text{age}_j = 0$ for the most recently written entry and
increases for older ones, and $\gamma \in (0, 1)$ is the decay factor.
Older entries contribute exponentially less to the cache logits,
implementing a form of soft forgetting without requiring explicit
eviction or additional bookkeeping.
When $\gamma = 1$ the method reduces to standard TDA.

\subsection{Stretch Goal --- Uncertainty-Aware Fusion}
\label{subsec:unc_method}

TDA's entropy criterion governs \emph{which samples enter} the cache.
A complementary question is: should the cache influence all test
images equally, regardless of how uncertain the model is about
the current image?
We argue that a highly uncertain image should trust the cached
prototypes less.
We therefore scale the effective adaptation weight at inference time:
\begin{equation}
  \alpha_\text{eff}(x) =
    \alpha \cdot \max\!\bigl(0,\; 1 - s \cdot H(x)\bigr),
\end{equation}
where $H(x) \in [0, 1]$ is the normalised entropy of the zero-shot
softmax prediction and $s \geq 0$ is a sensitivity parameter.
A confident image ($H \approx 0$) uses the full $\alpha$;
a near-uniform prediction ($H \approx 1$) gets negligible cache
support.
This is \emph{distinct} from TDA's eviction policy, which controls
what is stored in the cache.
Here we gate the cache's \emph{influence on the current prediction},
leaving the cache contents unchanged.

% =========================================================
\section{Experiments}
\label{sec:exp}

\subsection{Experimental Setup}

\paragraph{Datasets.}
We evaluate on two benchmarks following TDA~\cite{karmanov2024tda}.
The \emph{OOD benchmark} includes ImageNet-A~\cite{hendrycks2021imageneta}
(7,500 natural adversarial images, 200 classes) and
ImageNet-V2~\cite{recht2019imagenetv2}
(10,000 images, 1,000 classes, collected with a new pipeline).
The \emph{cross-domain benchmark} includes
DTD~\cite{cimpoi2014dtd} (1,692 texture images, 47 classes).
All evaluation is in the online single-image setting (batch size 1)
with no access to source training data, following the standard TTA
protocol.

\paragraph{Implementation.}
We use official pre-trained CLIP weights~\cite{radford2021clip} for
both CLIP-RN50 and CLIP-ViT-B/16, with the standard hand-crafted
class-name prompts.
TDA hyperparameters are set to their published values:
positive shot capacity $k{=}3$, negative $k{=}2$,
$\alpha{=}2.0$, $\beta{=}5.0$,
entropy thresholds $[\tau_l, \tau_h]{=}[0.2, 0.5]$,
negative mask threshold $p_l{=}0.03$.
Default values for our extensions are
$\tau_m{=}5.0$, $\mu{=}0.9$, $\gamma{=}0.999$, $s{=}1.0$.
All experiments run on a single NVIDIA RTX 5090 GPU under
\texttt{torch.no\_grad()}.
Latency is measured as mean per-image wall-clock time in milliseconds.

\subsection{Baseline Reproduction}
\label{subsec:repro}

Table~\ref{tab:repro} compares our re-implementation against the
published TDA numbers.
Our reproduced results are within $\pm0.4\%$ on all datasets for
both backbones, confirming that our codebase faithfully implements
the original algorithm.
Minor discrepancies are attributable to random seed and hardware
differences.
This shared experimental setup enables fair comparison: all our
control variants run under identical conditions and are compared
against the same reproduced baseline.
Full OOD baseline results across all five ImageNet variants are
reported in Appendix Table~\ref{tab:app-ood}.

\begin{table}[t]
\centering\small\setlength{\tabcolsep}{4pt}
\caption{Baseline reproduction. $\dagger$: published~\cite{karmanov2024tda}.
  X-Avg: mean top-1 accuracy over all 10 cross-domain datasets.}
\label{tab:repro}
\begin{tabular}{llccc}
\toprule
Method & Backbone & IN-A & IN-V2 & X-Avg \\
\midrule
CLIP$^\dagger$  & RN50     & 23.24 & 52.91 & 56.63 \\
TDA$^\dagger$   & RN50     & 30.29 & 55.54 & 61.03 \\
TDA (ours)      & RN50     & 30.98 & 55.19 & 61.23 \\
\midrule
CLIP$^\dagger$  & ViT-B/16 & 49.89 & 61.88 & 64.59 \\
TDA$^\dagger$   & ViT-B/16 & 60.11 & 64.67 & 67.53 \\
TDA (ours)      & ViT-B/16 & 60.43 & 64.37 & 67.91 \\
\bottomrule
\end{tabular}
\end{table}

\subsection{Control Ablations}
\label{subsec:controls}

Table~\ref{tab:controls} shows the effect of each control applied
in isolation and all three combined, using default hyperparameters.

On RN50, the margin filter causes a notable DTD accuracy drop from
44.09\% to 40.37\%.
The default threshold $\tau_m{=}5.0$ is calibrated for the logit
scale of ImageNet-1k, where the top-1 class typically dominates.
On DTD's 47-class fine-grained texture task, correct predictions
often have smaller margins, so many valid samples are filtered out.
The knob sweep in Section~\ref{subsec:ablations} shows that a
higher threshold ($\tau_m{=}10$--$20$) partly recovers this loss.

Momentum blending and affinity decay both produce near-neutral changes
when applied independently.
Their effects are bounded by the quality of the existing cache:
since TDA's entropy-based eviction already filters most incorrect
predictions, there is limited room for further improvement from
smoothing or aging.
Combining all three controls yields mixed results: a slight
improvement on ImageNet-A for ViT-B/16 margin (+0.02\%), but
a drop on DTD for both backbones.
Complete per-dataset accuracy for all 10 cross-domain datasets and
all control configurations is given in Appendix
Tables~\ref{tab:app-rn50} (RN50) and~\ref{tab:app-vitb16} (ViT-B/16).

\begin{table}[t]
\centering\small\setlength{\tabcolsep}{4pt}
\caption{Ablation of individual stability controls.
  Default hyperparameters throughout.}
\label{tab:controls}
\begin{tabular}{lcccc}
\toprule
Method & \multicolumn{2}{c}{RN50} & \multicolumn{2}{c}{ViT-B/16} \\
\cmidrule(lr){2-3}\cmidrule(lr){4-5}
 & IN-A & DTD & IN-A & DTD \\
\midrule
TDA baseline & 30.98 & 44.09 & 60.43 & \best{46.75} \\
+\,margin    & 30.81 & 40.37 & \best{60.45} & 44.44 \\
+\,momentum  & 30.85 & 43.62 & 60.17 & 45.45 \\
+\,decay     & 30.98 & 44.03 & 60.43 & 46.69 \\
all three    & 30.91 & 42.20 & 59.91 & 44.03 \\
\bottomrule
\end{tabular}
\end{table}

\subsection{Hyperparameter Knob Sweeps}
\label{subsec:ablations}

Table~\ref{tab:knobs} sweeps the primary hyperparameter of each
control on RN50, varying one at a time with the others held at default.

\paragraph{Margin threshold $\tau_m$.}
ImageNet-A accuracy is flat at 30.81\% for $\tau_m \leq 5$ and
rises slightly at larger values, reaching 31.06\% at $\tau_m{=}20$.
DTD shows a non-monotone pattern: it drops sharply at
$\tau_m{=}2$ (35.64\%) before recovering at higher thresholds.
The recovery at $\tau_m{=}10$--$20$ suggests that a stricter gate
has a net positive effect once it stops over-filtering, but the
optimal threshold is clearly dataset-dependent.

\paragraph{Momentum $\mu$.}
ImageNet-A accuracy increases monotonically with $\mu$, from 30.50\%
at $\mu{=}0$ (equivalent to always replacing the oldest entry) to
30.87\% at $\mu{=}0.95$.
DTD is maximised at $\mu{=}0.90$ (43.03\%).
Higher momentum stabilises the prototype by reducing the weight of
each new sample, which helps on a dataset as noisy as ImageNet-A
but can slow adaptation when the test distribution is homogeneous.

\paragraph{Decay factor $\gamma$.}
Mild decay at $\gamma{=}0.995$ gives the best ImageNet-A result
(31.03\%, +0.05\% over no decay).
The differences across the sweep are small in absolute terms,
indicating that the 7,500-image ImageNet-A stream is too short for
entries to accumulate significant age-related staleness.
Decay is expected to have a larger impact on longer streams or
streams with temporal correlation.

\begin{table}[t]
\centering\small\setlength{\tabcolsep}{2.8pt}
\caption{Hyperparameter knob sweeps on CLIP-RN50.
  Each column varies one parameter; others at default.
  \textbf{Bold}: best result per sweep.}
\label{tab:knobs}
\begin{tabular}{cc|cc|cc}
\toprule
$\tau_m$ & IN-A / DTD &
$\mu$    & IN-A / DTD &
$\gamma$ & IN-A / DTD \\
\midrule
1.0  & 30.81 / 41.19
     & 0.00  & 30.50 / 42.85
     & 1.000 & 30.98 / \best{43.56} \\
2.0  & 30.81 / 35.64
     & 0.50  & 30.60 / 42.85
     & 0.999 & 30.98 / 43.50 \\
5.0  & 30.81 / 40.37
     & 0.70  & 30.68 / 42.85
     & 0.995 & \best{31.03} / 43.44 \\
10.0 & 31.04 / 42.55
     & 0.90  & 30.85 / \best{43.03}
     & 0.990 & 31.01 / 43.38 \\
20.0 & \best{31.06} / \best{42.55}
     & 0.95  & \best{30.87} / 42.97
     & 0.950 & 31.00 / 43.20 \\
\bottomrule
\end{tabular}
\end{table}

\subsection{Shot Capacity and Latency Trade-off}
\label{subsec:capacity}

Table~\ref{tab:capacity} measures top-1 accuracy alongside
per-image wall-clock time as shot capacity $k$ varies.
Accuracy on ImageNet-A peaks at $k{=}5$ (31.00\%) but exceeds the
TDA default of $k{=}3$ by only 0.02\%, a negligible margin.
Latency grows sub-linearly: from 429\,ms at $k{=}1$ to 739\,ms
at $k{=}8$, a 72\% increase for a negligible accuracy gain.
Setting $k{=}1$ reduces latency by 33\% relative to $k{=}3$ while
losing only 0.13\% accuracy, making it an attractive operating point
for latency-sensitive deployments.

On DTD, accuracy peaks at $k{=}2$ (43.85\%) and degrades for
larger capacities.
This confirms TDA's own finding~\cite{karmanov2024tda}: on a small
47-class dataset, high shot capacity causes the cache to accumulate
noise, and the accuracy drop at $k{=}8$ (42.91\%) is principally
due to noisy pseudo-labels rather than class imbalance.

\begin{table}[t]
\centering\small\setlength{\tabcolsep}{5pt}
\caption{Shot capacity vs.\ accuracy and latency, CLIP-RN50.
  Latency is mean per-image wall-clock time on ImageNet-A.}
\label{tab:capacity}
\begin{tabular}{cccc}
\toprule
$k$ & IN-A (\%) & DTD (\%) & Latency (ms) \\
\midrule
1 & 30.85 & 43.79          & \best{429} \\
2 & 30.91 & \best{43.85}   & 664 \\
3 & 30.98 & 43.56          & 640 \\
5 & \best{31.00} & 43.85   & 720 \\
8 & 30.84 & 42.91          & 739 \\
\bottomrule
\end{tabular}
\end{table}

\subsection{Uncertainty-Aware Adaptation}
\label{subsec:unc}

Table~\ref{tab:unc} sweeps the sensitivity parameter $s$ for both
backbones, with $s{=}0$ corresponding to the unmodified TDA baseline.

For RN50, $s{=}0.7$ gives the best DTD accuracy (43.85\%,
+0.29\% over baseline), while ImageNet-A shows a small monotonic
decrease with increasing $s$.
The interpretation is that on DTD, RN50 frequently makes uncertain
predictions, and down-weighting their cache contribution reduces the
impact of noisy updates.
For ViT-B/16, all values of $s{>}0$ reduce accuracy on both
datasets.
ViT-B/16 is a much stronger encoder; its zero-shot confidence is
well-calibrated on these benchmarks, so filtering out uncertain images
removes genuinely informative samples rather than noise.

Combining all three controls with uncertainty scaling (last row of
Table~\ref{tab:unc}) does not improve over the plain TDA baseline
for either backbone.
This indicates that the controls interact: the margin filter and
uncertainty scaling both suppress low-confidence samples, leading
to double-filtering and over-restriction of cache updates.

\begin{table}[t]
\centering\small\setlength{\tabcolsep}{4pt}
\caption{Uncertainty-aware adaptation strength sweep.
  $s{=}0$: TDA baseline.
  ``All'': margin + momentum + decay combined, $s{=}0.7$.}
\label{tab:unc}
\begin{tabular}{lcccc}
\toprule
$s$ & \multicolumn{2}{c}{RN50} & \multicolumn{2}{c}{ViT-B/16} \\
\cmidrule(lr){2-3}\cmidrule(lr){4-5}
 & IN-A & DTD & IN-A & DTD \\
\midrule
0 (baseline) & 30.98 & 43.56 & 60.43 & \best{46.75} \\
0.3          & 30.97 & 43.62 & 60.38 & 44.98 \\
0.5          & 30.97 & 43.62 & 60.32 & 44.86 \\
0.7          & 30.94 & \best{43.85} & 60.29 & 44.80 \\
1.0          & 30.91 & 43.44 & 60.27 & 44.68 \\
\midrule
All + unc    & 30.88 & 42.43 & 59.94 & 44.27 \\
\bottomrule
\end{tabular}
\end{table}

\subsection{Discussion}
\label{subsec:disc}

\paragraph{Why are accuracy improvements modest?}
The overall accuracy changes across all controls are small
(below 1\%).
For ImageNet-A with RN50, zero-shot CLIP accuracy is only
$\sim\!23\%$, suggesting that $\sim\!77\%$ of raw predictions are
wrong.
However, TDA's entropy-based eviction policy already achieves
$\sim\!90\%$ cache correctness because wrong predictions tend to have
high entropy and are filtered out~\cite{karmanov2024tda}.
Our margin filter provides an orthogonal signal but on an
already-high-quality cache, leaving limited headroom for improvement.

\paragraph{Latency overhead.}
All three controls add negligible compute overhead compared to the
CLIP forward pass ($\sim$410\,ms for RN50, $\sim$620\,ms for
ViT-B/16 on our hardware).
The margin filter requires only a two-value subtraction from
already-computed logits.
Momentum blending adds one vector interpolation and normalisation
call (${\sim}2\,\mu$s for $d{=}1{,}024$).
Affinity decay is a scalar elementwise multiply applied once per
retrieval.
None of the controls require additional forward passes or stored
activations beyond what TDA already maintains.

% =========================================================
\section{Conclusion}

We reproduced TDA within $\pm0.4\%$ of published results on OOD and
cross-domain benchmarks with two CLIP backbones, providing a verified
baseline for fair comparison.
We proposed and ablated four streaming stability controls: a margin
filter, momentum blending, affinity decay, and uncertainty-aware
adaptation strength, each independently togglable without code changes.
Individual controls have modest but interpretable effects: the margin
filter is dataset-dependent and requires tuning to avoid
over-filtering fine-grained datasets; momentum and decay both benefit
from higher values but saturate quickly; uncertainty scaling helps the
weaker RN50 encoder but hurts the well-calibrated ViT-B/16.
Combining all controls does not consistently improve over the baseline,
revealing that independently designed filters can interact and
double-filter samples.
The most actionable finding is the accuracy--latency trade-off from
the shot capacity sweep: $k{=}1$ reduces per-image latency by 33\%
at only $-0.13\%$ accuracy cost, offering a practical configuration
for real-time applications.
Future work should investigate adaptive thresholds that adjust to the
current stream's difficulty, and evaluate stability controls under
class-imbalanced or temporally correlated test streams.

% =========================================================
% References --- do not count toward the page limit
{\small\bibliographystyle{unsrt}\bibliography{references}}

% =========================================================
% APPENDIX --- does not count toward the 5-page limit
% =========================================================
\clearpage
\onecolumn
\appendix
\section*{Appendix: Full Results Tables}

\noindent Tables~\ref{tab:app-ood}, \ref{tab:app-rn50}, and
\ref{tab:app-vitb16} report complete per-dataset accuracy (\%) for
both backbones across every dataset used in the study.  All runs use
the same fixed hyperparameters as the main paper
($\tau_m{=}5.0$, $\mu{=}0.9$, $\gamma{=}0.999$, shot capacity~3).

% ------------------------------------------------------------------
\subsection*{A.\; OOD Benchmark --- Baseline Performance}

\begin{table}[h]
\centering
\caption{Baseline (vanilla TDA, no extra controls) accuracy on the
  full OOD benchmark.  \emph{IN} = ImageNet, \emph{IN-A} =
  ImageNet-A, \emph{IN-V2} = ImageNet-V2, \emph{IN-R} = ImageNet-R,
  \emph{IN-S} = ImageNet-Sketch.}
\label{tab:app-ood}
\setlength{\tabcolsep}{8pt}
\begin{tabular}{lccccc}
\toprule
Backbone & IN & IN-A & IN-V2 & IN-R & IN-S \\
\midrule
RN50        & 61.31 & 30.98 & 55.19 & 62.80 & 37.97 \\
ViT-B/16    & 69.58 & 60.43 & 64.37 & 80.62 & 51.43 \\
\bottomrule
\end{tabular}
\end{table}

% ------------------------------------------------------------------
\subsection*{B.\; Cross-Domain Benchmark --- RN50}

\begin{table}[h]
\centering
\caption{Per-dataset accuracy (\%) on the 10-dataset cross-domain
  benchmark for the RN50 backbone under each stability-control
  configuration.  \emph{Avg} is the unweighted mean over all 10
  datasets.}
\label{tab:app-rn50}
\setlength{\tabcolsep}{4.5pt}
\small
\begin{tabular}{lccccccccccr}
\toprule
Control
  & \rotatebox{60}{Caltech} & \rotatebox{60}{DTD}
  & \rotatebox{60}{EuroSAT} & \rotatebox{60}{FGVC}
  & \rotatebox{60}{Food} & \rotatebox{60}{Flowers}
  & \rotatebox{60}{Pets} & \rotatebox{60}{Cars}
  & \rotatebox{60}{SUN} & \rotatebox{60}{UCF}
  & \rotatebox{60}{Avg} \\
\midrule
Baseline       & 89.37 & 44.09 & 45.84 & 16.41 & 77.65 & 68.62 & 85.94 & 57.49 & 62.74 & 64.16 & 61.23 \\
+Margin        & 89.09 & 40.37 & 35.49 & 12.54 & 77.60 & 67.32 & 85.53 & 56.40 & 62.31 & 63.76 & 59.04 \\
+Momentum      & 88.40 & 43.62 & 38.96 & 16.71 & 77.75 & 68.09 & 86.29 & 56.67 & 62.34 & 64.39 & 60.32 \\
+Decay         & 89.37 & 44.03 & 45.77 & 16.44 & 77.64 & 68.57 & 86.02 & 57.46 & 62.74 & 64.16 & 61.22 \\
All three      & 87.42 & 42.20 & 36.95 & 16.29 & 77.27 & 67.80 & 85.09 & 56.49 & 60.82 & 63.52 & 59.39 \\
\bottomrule
\end{tabular}
\end{table}

% ------------------------------------------------------------------
\subsection*{C.\; Cross-Domain Benchmark --- ViT-B/16}

\begin{table}[h]
\centering
\caption{Per-dataset accuracy (\%) on the 10-dataset cross-domain
  benchmark for the ViT-B/16 backbone under each stability-control
  configuration.  \emph{Avg} is the unweighted mean over all 10
  datasets.}
\label{tab:app-vitb16}
\setlength{\tabcolsep}{4.5pt}
\small
\begin{tabular}{lccccccccccr}
\toprule
Control
  & \rotatebox{60}{Caltech} & \rotatebox{60}{DTD}
  & \rotatebox{60}{EuroSAT} & \rotatebox{60}{FGVC}
  & \rotatebox{60}{Food} & \rotatebox{60}{Flowers}
  & \rotatebox{60}{Pets} & \rotatebox{60}{Cars}
  & \rotatebox{60}{SUN} & \rotatebox{60}{UCF}
  & \rotatebox{60}{Avg} \\
\midrule
Baseline       & 94.04 & 46.75 & 60.22 & 24.45 & 86.05 & 71.70 & 89.94 & 67.19 & 67.89 & 70.82 & 67.91 \\
+Margin        & 93.51 & 44.44 & 56.04 & 18.87 & 85.99 & 70.56 & 89.42 & 65.80 & 67.54 & 68.09 & 66.03 \\
+Momentum      & 94.04 & 45.45 & 49.67 & 24.99 & 86.03 & 70.69 & 89.59 & 67.03 & 67.71 & 68.81 & 66.40 \\
+Decay         & 94.04 & 46.69 & 60.21 & 24.45 & 86.04 & 71.70 & 90.00 & 67.17 & 67.90 & 70.82 & 67.90 \\
All three      & 93.27 & 44.03 & 40.40 & 25.08 & 85.55 & 69.31 & 87.38 & 66.51 & 65.69 & 68.44 & 64.57 \\
\bottomrule
\end{tabular}
\end{table}

\end{document}
